% Source: https://artofproblemsolving.com/community/c3h2646133_formula_for_sum_of_divisors_of_number
% Stable identity: aops-topic-2646133-q01
Every divisor of $n$ is $p_1^{e_1}\cdots p_k^{e_k}$ with $0\le e_i\le a_i$, and by unique factorisation each divisor arises from exactly one choice of exponents. Expanding the product
\[\left(1+p_1+\cdots+p_1^{a_1}\right)\left(1+p_2+\cdots+p_2^{a_2}\right)\cdots\left(1+p_k+\cdots+p_k^{a_k}\right)\]
gives one term $p_1^{e_1}\cdots p_k^{e_k}$ for each choice of exponents, i.e. each divisor exactly once. Summing each geometric series,
\[\boxed{\sigma(n)=\prod_{i=1}^{k}\dfrac{p_i^{a_i+1}-1}{p_i-1}}.\]
For $360=2^3\cdot3^2\cdot5$: $\sigma(360)=(1+2+4+8)(1+3+9)(1+5)=15\cdot13\cdot6=\boxed{1170}$.
